% Consulta oficial a los responsables de los datos (formulario PDF editable).
% Autónomo: no depende del release ni del preámbulo de los manuscritos.
\documentclass[11pt]{article}
\usepackage[T1]{fontenc}
\usepackage[utf8]{inputenc}
\usepackage{stix2}
\usepackage[spanish,es-noshorthands,es-nolists]{babel}
\usepackage[a4paper,margin=2.3cm]{geometry}
\usepackage{microtype}
\usepackage{xcolor}
\usepackage{array}
\usepackage{booktabs}
\usepackage{tabularx}
\usepackage{needspace}
\usepackage{enumitem}
\usepackage[hidelinks,pdfusetitle]{hyperref}
\hypersetup{pdftitle={Consulta a los responsables de los datos: VA-MENGOC-BC},
  pdfauthor={Richard Alejandro Matos Arderi; Abel Ponce Gonzalez},
  pdfsubject={Formulario editable de preguntas abiertas sobre las bases de datos de VA-MENGOC-BC}}

\definecolor{azul}{HTML}{1C5CAB}
\definecolor{gris}{HTML}{8A8A85}
\definecolor{fondo}{HTML}{FFFFF0}
\setlength{\parindent}{0pt}
\setlength{\parskip}{3pt}
\emergencystretch=3em
\setlist{nosep}

\newcounter{preg}
\newcommand{\Fichero}[1]{\textit{«#1»}}
% Pregunta: id, título, referencia, texto, tratamiento actual, alto del cuadro (cm)
\newcommand{\Pregunta}[6]{%
  \Needspace{\dimexpr 9\baselineskip+#6cm\relax}%
  \refstepcounter{preg}%
  {\color{azul}\textbf{Pregunta \thepreg.}}\ \textbf{#2}\par
  \textit{Referencia:} #3\par
  #4\par
  \textit{Tratamiento actual en el análisis:} #5\par
  \smallskip
  \TextField[name=q\thepreg#1,multiline=true,width=\linewidth,height=#6cm,charsize=10pt,%
    bordercolor=gris,backgroundcolor=fondo,align=0]{}\par
  \bigskip}
\newcommand{\Seccion}[1]{\Needspace{14\baselineskip}\section*{\color{azul}#1}}
\newcommand{\Campo}[2]{\noindent\makebox[5.4cm][l]{\textbf{#1}}\TextField[name=#2,%
  width=\dimexpr\linewidth-5.6cm\relax,height=0.6cm,charsize=10pt,bordercolor=gris,backgroundcolor=fondo]{}\par\smallskip}

\begin{document}
\begin{Form}
\raggedright
\begin{center}
{\small\textcolor{azul}{\textbf{DOCUMENTO DE CONSULTA}}\par}
\medskip
{\Large\bfseries Consulta a los responsables de los datos\par}
\medskip
{\large Proyecto VA-MENGOC-BC: farmacovigilancia nacional y calidad de lotes\par}
\medskip
{\small Dirección de Investigaciones Clínicas y Evaluación de Impacto (DICEI)\\
Instituto Finlay de Vacunas \textperiodcentered{} 16 de noviembre de 2026\par}
\end{center}

\section*{\color{azul}Propósito y modo de respuesta}
El grupo de ciencia de datos de la DICEI prepara dos artículos científicos sobre la vacuna
VA-MENGOC-BC: el primero, sobre la vigilancia poscomercialización de los eventos supuestamente
atribuibles a la vacunación o inmunización (ESAVI) en Cuba; el segundo, sobre la relación entre
los atributos de calidad de los lotes liberados y los eventos notificados. El análisis se ejecuta
en la estación controlada del Instituto con las bases de datos originales; ningún dato individual
sale de ella.

Algunas decisiones metodológicas dependen de información que solo conocen quienes administran,
registran o custodian las bases de datos. Se solicita su respuesta a las preguntas siguientes.
Cada una indica el fichero, la hoja y, cuando procede, la columna a la que se refiere, y el
tratamiento que el análisis aplica de forma provisional, de modo que baste confirmarlo o
corregirlo. Este documento no contiene datos individuales: solo recuentos agregados.

\textbf{Cómo responder.} Haga clic en cada recuadro amarillo y escriba; los recuadros no tienen
límite de extensión (si el texto excede el recuadro, este se desplaza). Puede responder solo las
preguntas de su competencia e indicar en el resto «no es mi competencia» o el nombre de quien
puede responder. Abra el documento en un navegador (Chrome, Edge o Firefox) o en un lector de PDF,
guárdelo con el nombre «consulta\_respondida» (Ctrl+S o «Guardar como») y devuélvalo al grupo de
ciencia de datos de la DICEI. Las respuestas se incorporan a la configuración versionada del
proyecto y se registran con su fecha y responsable.

\subsection*{Datos de quien responde}
\Campo{Nombre y apellidos}{respnombre}
\Campo{Cargo}{respcargo}
\Campo{Institución o dependencia}{respinstitucion}
\Campo{Fecha de respuesta}{respfecha}
\Campo{Preguntas que contesta}{resppreguntas}

\Seccion{A. Procedencia, autorización y custodia de los datos}
\Pregunta{extraccion}{Fecha de extracción y custodio de la base de datos}%
{Libros anuales de ESAVI \Fichero{EA 2017 TOTAL.xlsx} a \Fichero{EA 2025 TOTAL.xlsx} y los libros de 2026.}%
{¿En qué fecha se extrajo cada libro del sistema de origen, o desde qué fecha de corte son válidos?
¿Qué institución o dependencia custodia la base de datos nacional de ESAVI y debe figurar como
tal en los artículos?}%
{Ambos datos se imprimen en los manuscritos como «a completar por los autores» hasta recibir su respuesta.}{3.4}

\Pregunta{etica}{Autorización y base legal para el uso de los datos}%
{Todos los ficheros de ESAVI y el libro de lotes \Fichero{BD para publicación lotes final.xlsx}.}%
{¿Existe aprobación de un comité de ética de la investigación o autorización de la dirección
institucional para el uso secundario de estos datos? Indique el órgano, el número y la fecha de la
resolución, y la base legal aplicada (Ley 149/2022 de protección de datos personales u otra).}%
{Tratamiento local con seudonimización irreversible de nombres y direcciones (resumen con clave
secreta), fecha de nacimiento convertida en edad y publicación solo de agregados con supresión de
recuentos menores que 5.}{3.4}

\Pregunta{retencion}{Retención de los datos y evaluación de impacto en la privacidad}%
{Capas de datos restringidas (\Fichero{data/raw}, \Fichero{data/curated}) de la estación controlada.}%
{¿Qué plazo de conservación establece la política institucional para los ficheros originales y para
la capa seudonimizada? ¿Se exige una evaluación de impacto en la protección de datos antes de la
publicación de resultados agregados, y quién la realiza?}%
{Los ficheros originales son de solo lectura, nunca se versionan ni se envían a servicios externos;
no se ha fijado un plazo de retención.}{3.2}

\Pregunta{lotesabs}{Atributos de calidad publicables en unidades absolutas}%
{\Fichero{BD para publicación lotes final.xlsx}, hoja 1 (control de calidad de lotes), columnas F a N:
pH, concentraciones y adsorciones de proteína y polisacárido, Al(OH)\textsubscript{3}, tiomersal, inmunogenicidad
(actividad bactericida e IgG) y endotoxinas.}%
{Los datos de lotes son información industrial confidencial. ¿Cuáles de estos atributos pueden
publicarse en unidades absolutas (por ejemplo, en tablas y figuras de un artículo) y cuáles solo
como posición dentro de la ventana de especificación? ¿Pueden publicarse los límites de
especificación?}%
{Solo se publica la posición normalizada dentro de la ventana de especificación y los índices de
capacidad; ningún valor absoluto ni límite sale de la estación.}{3.4}

\Seccion{B. Libro de lotes \Fichero{BD para publicación lotes final.xlsx}}
\Pregunta{inocuidad}{Significado de «-» en la prueba de inocuidad}%
{Hoja 1 (control de calidad de lotes), columna «Inocuidad»: 108 de las 737 filas contienen «-»; el
resto es «Satisfactoria».}%
{¿El guion indica que el ensayo no se realizó, que no era aplicable o que el campo no era
obligatorio en determinados períodos? Si es así, ¿cuáles?}%
{Se interpreta como «no determinado»; no se cuenta como ensayo no satisfactorio.}{3.0}

\Pregunta{filasdup}{Filas repetidas por lote}%
{Hoja 1, columna «Lote de vacuna envasada»: 737 filas para 714 lotes distintos (22 identificadores
aparecen en más de una fila).}%
{¿Una misma identificación de lote puede corresponder legítimamente a varias filas (varios destinos
o varios ensayos del mismo lote) o se trata de lotes distintos que reutilizan la identificación, o de
errores de registro? ¿Qué fila debe prevalecer si los valores difieren?}%
{Se mantiene cada fila con una clave compuesta (lote, año, orden) y, al enlazar una notificación,
se prefiere el lote nacional cuya vigencia cubre la fecha de vacunación.}{3.4}

\Pregunta{destino}{Destino de los lotes}%
{Hoja 1, columna «Destino»: 21 categorías (Vietnam, Nacional, Venezuela, Colombia, entre otras) y un
valor «vietnam colombia».}%
{¿«Nacional» designa los lotes destinados al programa nacional de inmunización? ¿El valor
«vietnam colombia» representa un lote repartido entre dos mercados o un error de registro? ¿Algún
destino debe agruparse con otro?}%
{«Nacional» frente a exportación como único criterio de agrupación; los valores se normalizan por
mayúsculas, acentos y espacios.}{3.0}

\Pregunta{codigo}{Codificación del número de lote}%
{Hoja 1, columna «Lote de vacuna envasada»: tres formatos, 999A (71,1 \%), A999 (23,4 \%) y A999A
(2,7 \%); columna «Año producción de la vacuna» (2011--2025).}%
{El primer dígito del identificador es la última cifra del año de liberación del lote (un lote de
2001 o de 2011 empieza por 1). ¿Esta regla se aplica por igual a los tres formatos, incluido el
formato con letra inicial (A999)? ¿Qué codifican las letras? ¿El «año de producción» de la columna C
coincide siempre con el año de liberación o puede diferir en un año?}%
{La regla se aplica al primer dígito del identificador en todos los formatos y se comprueba contra
la columna C: debe reproducirla (o el año siguiente) en al menos el 90 \% de los lotes.}{3.6}

\Pregunta{vigencia}{Cómputo de la vigencia de tres años}%
{Mismo libro y columnas que la pregunta anterior; notificaciones de ESAVI de los libros anuales.}%
{La vigencia de un lote es de tres años. ¿Se cuenta desde la fecha de liberación, desde la de
producción o hasta el fin del año natural correspondiente? ¿Existen prórrogas de vigencia
autorizadas para algún lote, o lotes con una vigencia distinta?}%
{Un lote se considera vencido si se usa después del 31 de diciembre del año de liberación más tres.
Las notificaciones con un lote vencido o aún no liberado se excluyen del análisis por lote.}{3.4}

\Seccion{C. Libros anuales de ESAVI (\Fichero{EA 2017 TOTAL.xlsx} a \Fichero{EA 2025 TOTAL.xlsx})}
\Pregunta{duplicados}{Filas exactamente repetidas}%
{Hoja principal de cada libro; por ejemplo, 119 filas duplicadas en \Fichero{EA 2017 TOTAL.xlsx}, 18 en
\Fichero{EA 2018 TOTAL.xlsx}, 19 en \Fichero{EA 2019 TOTAL.xlsx}, 9 en \Fichero{EA 2020 TOTAL.xlsx},
170 en \Fichero{EA 2024 TOTAL.xlsx} y 3 en \Fichero{EA 2025 TOTAL.xlsx}.}%
{¿Las filas idénticas en todas sus columnas son transcripciones repetidas de una misma notificación o
pueden corresponder a vacunaciones distintas con los mismos datos (por ejemplo, dos personas con
iguales datos o dos dosis el mismo día)?}%
{Se elimina la repetición exacta conservando la primera aparición; además se fusionan los registros
con la misma persona seudonimizada, fecha de vacunación, vacunas y lotes.}{3.4}

\Pregunta{hoja2018}{Hoja de casos que no entraron en el cierre de 2018}%
{\Fichero{EA 2018 TOTAL.xlsx}, hoja «Casos que no entraronen el cier» (33 filas, octubre a diciembre
de 2018, con la misma estructura que «TOTAL EA»).}%
{¿Estos casos son notificaciones adicionales que no figuran en la hoja «TOTAL EA» y deben sumarse al
análisis de 2018, o son copia de casos ya incluidos? ¿Existen hojas equivalentes en otros años?}%
{Se leen además de la hoja principal y se eliminan las repeticiones exactas con ella.}{3.0}

\Pregunta{ea2023}{Cobertura temporal del libro de 2023}%
{\Fichero{EA 2023 TOTAL.xlsx}, hoja «EA TOTAL» (5072 filas): la columna MES solo contiene los meses 1 a
7.}%
{¿Existen notificaciones de agosto a diciembre de 2023? En caso afirmativo, ¿en qué fichero o sistema
están? Si no existen, ¿se debe a que no se recogieron, a que no se consolidaron o a que no hubo casos?}%
{Los meses sin datos se tratan como ausentes (no como cero) y 2023 se considera año parcial, sin tasa
anual de notificación.}{3.2}

\Pregunta{fechas}{Fechas imposibles}%
{\Fichero{EA 2019 TOTAL.xlsx}, hoja «Hoja1»: F.VACUNACION con años hasta 2029 y F.NOTIFICACION desde 2002;
\Fichero{EA 2017 TOTAL.xlsx} y \Fichero{EA 2020 TOTAL.xlsx}: F.NACIMIENTO con años como 1699, 1905,
2075 y 2107.}%
{¿Existe una fuente primaria (ficha de notificación u otra) con la que corregir estas fechas? ¿Considera adecuada la regla de corrección indicada
en el tratamiento actual? ¿Hay otro error de tecleo frecuente (por ejemplo, día y mes invertidos)
que deba tratarse?}%
{El año de vacunación o de notificación se corrige solo cuando el día y el mes, la otra fecha y el
año del libro lo determinan con una demora de 0 a 365 días; si no, la fecha no se usa. Las fechas de
nacimiento imposibles no se corrigen y se usa la edad registrada.}{3.6}

\Pregunta{mes}{Significado de la columna MES}%
{Todos los libros anuales, columna MES (valores 1 a 12).}%
{¿MES es el mes de la vacunación, el de la notificación o el de la ocurrencia del evento? ¿Puede
diferir de F.VACUNACION?}%
{Se usa como control de calidad y para determinar qué meses cubre cada libro; no define el año de
análisis, que es el de la fecha de vacunación.}{2.8}

\Pregunta{colfinal}{Columna final sin encabezado}%
{\Fichero{EA 2017 TOTAL.xlsx}, hoja «TOTAL», columna BI (a continuación de SECUELA), sin título y casi
vacía.}%
{¿Qué representa esta columna? ¿Debe incorporarse al análisis o descartarse?}%
{Se descarta y se informa como columna sin reconocer.}{2.6}

\Pregunta{covid}{Hojas por producto de vacunas contra la COVID-19}%
{\Fichero{EA 2021 TOTAL.xlsx}, \Fichero{EA 2022 TOTAL.xlsx}, \Fichero{EA 2023 TOTAL.xlsx} y
\Fichero{EA 2024 TOTAL.xlsx}: hojas como ABDALA, SOBERANA y SOBERANA Plus, además de las hojas de
intervenciones por municipio.}%
{¿Estas hojas son extractos de la hoja principal o contienen notificaciones que no figuran en ella? En
particular, ¿la hoja principal incluye todos los casos de SOBERANA (producto del Instituto Finlay) y de
ABDALA?}%
{Las hojas por producto no se leen (se consideran copias filtradas de la hoja principal); los casos
de la hoja principal con estos productos se clasifican como SOBERANA (Finlay) y ABDALA (CIGB).}{3.4}

\Seccion{D. Formulario de notificación utilizado desde 2024 (\Fichero{EA 2024 TOTAL.xlsx} y \Fichero{EA 2025 TOTAL.xlsx})}
\Pregunta{fiebre}{Equivalencia de FIEBRE con los indicadores anteriores}%
{Columna FIEBRE de \Fichero{EA 2024 TOTAL.xlsx} y \Fichero{EA 2025 TOTAL.xlsx}, junto a FEBRICULA; hasta
2023 existían FIEBRE 39 y FIEBRE 40.}%
{¿FIEBRE equivale a fiebre de 39\,°C o más, a 40\,°C o más, o a una definición distinta (por ejemplo,
cualquier temperatura mayor que la de FEBRICULA)? ¿Qué umbral de temperatura define FEBRICULA?}%
{Cada indicador se trata como distinto y se analiza solo en los años cuyo formulario lo contiene; un
indicador ausente de un año es desconocido, no «no notificado».}{3.4}

\Pregunta{local}{Reacciones locales: columna combinada y absceso}%
{Columna «DOLOR, ENROJECIMIENTO, INDURACION» y columna ABSCESO de los libros de 2024 y 2025; hasta
2023 existían R.L.S., ABS.ESTERIL y ABS.BACTERIANA.}%
{¿La columna combinada se marca «S» si ocurre cualquiera de las tres manifestaciones? ¿Equivale a la
reacción local severa (R.L.S.) de los años anteriores o es una categoría más amplia? ¿ABSCESO
corresponde a absceso estéril, bacteriano o ambos? ¿R.L.S. dejó de registrarse?}%
{La columna combinada y ABSCESO son indicadores nuevos, sin equivalencia asumida con R.L.S. ni con los
abscesos de años anteriores.}{3.6}

\Pregunta{otrosind}{Otros indicadores nuevos o desaparecidos}%
{Libros de 2024 y 2025: ANAFILAXIA, TROMBOCITOPENIA, EHH, TA ALTA, DECAIMIENTO, CEFALEA, CRISIS VAGAL,
ERITEMA, SOMNOLENCIA, EDEMA LABIAL, EDEMA FACIAL y PARALISIS FACIAL; ya no figuran LLANTO
PERSISTENTE, COLAPSO, C.FEBRIL, C.AFEBRIL, PURPURA, R.ANAFILACTICA y otros de 2017 a 2023.}%
{¿ANAFILAXIA equivale a R.ANAFILACTICA y TROMBOCITOPENIA a PURPURA? ¿Qué significa EHH (¿episodio
hipotónico-hiporreactivo?)? ¿Los indicadores desaparecidos dejaron de registrarse o cambiaron de
nombre? ¿Dónde se anotan ahora el llanto persistente o las convulsiones?}%
{Sin equivalencias asumidas; cada indicador se analiza solo en los años que lo registran.}{3.8}

\Pregunta{app}{Antecedentes patológicos personales (APP)}%
{\Fichero{EA 2024 TOTAL.xlsx}: columnas APP ALERGIA, APP CONVULSON, APP ASMA, APP HTA, APP DM y APP CI;
\Fichero{EA 2025 TOTAL.xlsx}: columna APP (5 categorías) además de las anteriores.}%
{¿Qué significan APP CI (¿cardiopatía isquémica?) y la columna APP del libro de 2025? ¿Cómo se relacionan
con los antecedentes de alergia, convulsión y asma de 2017 a 2023, que incluían antecedentes
familiares?}%
{Los antecedentes se conservan como indicadores propios del formulario de cada año; APP CI se trata
como dato sensible y no se publica.}{3.2}

\Pregunta{grave}{Clasificación GRAVE y NO GRAVE}%
{\Fichero{EA 2024 TOTAL.xlsx} y \Fichero{EA 2025 TOTAL.xlsx}, columnas GRAVE y NO GRAVE (esta última con 10,4 \%
de valores vacíos en 2024).}%
{¿Qué criterio define «grave» (definición de la OMS u otra)? ¿Las dos columnas son complementarias?
¿Cómo debe interpretarse un valor vacío? Además, ¿existe una clasificación de gravedad por evento y
alguna codificación MedDRA o CIE-10 en el sistema de origen?}%
{Se usa la clasificación de la OMS (común, rara, grave) y las definiciones de la Brighton Collaboration
cuando existen; las columnas GRAVE y NO GRAVE se conservan como dato notificado.}{3.6}

\Pregunta{carne}{Uso del carné de identidad}%
{\Fichero{EA 2024 TOTAL.xlsx} y \Fichero{EA 2025 TOTAL.xlsx}, columna CARNE DE IDENTIDAD (13,1 \% de valores
vacíos en 2024).}%
{¿Se autoriza su uso, dentro de la estación controlada, como clave adicional para detectar
notificaciones repetidas de una misma persona, o debe descartarse en todos los casos?}%
{Se reconoce como identificador directo y se descarta al ingresar; el seudónimo se construye con
nombre, fecha de nacimiento y sexo en todos los años.}{3.0}

\Seccion{E. Esquema de vacunación, siglas de vacunas y comparador}
\Pregunta{esquema}{Esquema de vacunación de VA-MENGOC-BC}%
{Columnas EDAD, TIPO DE VACUNA (AM-BC) y DOSIS de todos los libros anuales.}%
{Se entiende que la dosis 1 se administra a los 3 meses y la dosis 2 a los 5 meses de edad, con un
intervalo mínimo de 6 a 8 semanas y sin refuerzo en el uso programado. ¿Es correcto para todo el
período 2017--2026? ¿Hubo campañas de recuperación o usos fuera del esquema (niños mayores o
adultos, refuerzo «R», dosis 3 o 4) y en qué años? ¿Qué edad máxima corresponde a la indicación?}%
{Son válidas las dosis 1 y 2 desde 2,87 y 4,25 meses de edad (con 4 días de gracia); se excluyen por
error de registro las dosis por debajo de esa edad y las dosis 3 o superiores; los refuerzos «R» y
las dosis tardías de recuperación se conservan.}{4.0}

\Pregunta{dosiscod}{Notación de la dosis con dos vacunas}%
{Columna DOSIS de todos los libros: valores como «2.1», «2 y 1», «1,1» y «2,1».}%
{¿Estas notaciones significan una dosis por cada vacuna administrada en la misma sesión (por ejemplo,
dosis 2 de una vacuna y dosis 1 de otra), en el orden en que aparecen en TIPO DE VACUNA?}%
{Se interpretan así, en el mismo orden que las vacunas de la columna TIPO DE VACUNA.}{3.0}

\Pregunta{siglas}{Siglas de la columna TIPO DE VACUNA}%
{Todos los libros anuales, columna TIPO DE VACUNA. Significados asumidos: AM-BC, VA-MENGOC-BC;
PENTA-L, Heberpenta-L (CIGB); DPT, DT y TT, vacunas del Instituto Finlay; AT, antitifoídica Vi; AL,
antileptospirósica; PRS, triple viral; AG, antigripal; HB, hepatitis B; HIB, Hib conjugada; OPV e IPV,
poliomielitis oral e inactivada; PCV-10, Pneumosil-10 (Serum Institute of India, esquema desde 2024).}%
{Corrija o complete esta lista. En particular, ¿qué significa AA? ¿Se confirman AG y AL? ¿AT designa la
antitifoídica o la antitetánica? ¿PCV-10 es siempre Pneumosil-10? ¿Qué otras siglas aparecen y a qué
vacuna y fabricante corresponden?}%
{Mapeo marcado «inferido» en el diccionario de vacunas; AA queda sin significado confirmado.}{4.4}

\Pregunta{provincias}{Códigos de provincia}%
{Columna PROVINCIA o PROVINCIAS de todos los libros anuales (16 códigos).}%
{Se interpretan PR, AT, LH, MY, MT, IJ, VC, CF, SS, CA, CM, LT, HO, GM, SC y GT como Pinar del Río,
Artemisa, La Habana, Mayabeque, Matanzas, Isla de la Juventud, Villa Clara, Cienfuegos, Sancti
Spíritus, Ciego de Ávila, Camagüey, Las Tunas, Holguín, Granma, Santiago de Cuba y Guantánamo.
¿Es correcta la asignación, en especial GM (¿Granma?) y CH?}%
{Esa asignación; el código CH se agrupa con La Habana.}{3.0}

\Pregunta{comparador}{Comparador del primer artículo: productos del Instituto Finlay}%
{Columnas TIPO DE VACUNA y FABRICANTE de todos los libros anuales.}%
{Se acordó con la jefatura del proyecto que VA-MENGOC-BC se compare solo con las vacunas propias del
Instituto Finlay y con los tratamientos concomitantes. ¿Qué productos de TIPO DE VACUNA deben
considerarse del Instituto Finlay en cada período? Se han supuesto DPT, DT, TT, antitifoídica Vi (AT),
antileptospirósica (AL), Quimi-Vio y SOBERANA. ¿Los valores FINLAY, INST FINLAY, INSTITUTO FINLAY e IFV de
FABRICANTE designan siempre a este Instituto? ¿Qué otro producto debe incluirse o excluirse?}%
{Comparadores: los productos anteriores; las vacunas de otros fabricantes solo se analizan como
concomitantes de VA-MENGOC-BC.}{4.0}

\Pregunta{concomitantes}{Tratamientos concomitantes}%
{Todos los libros anuales (columnas TIPO DE VACUNA, DOSIS, S.APLICACION, V.ADMON., LOTE y CUAL (otro)).}%
{¿Alguna columna registra tratamientos concomitantes distintos de las vacunas (por ejemplo, antipiréticos
o antibióticos) o el uso de otros productos el mismo día? ¿Cómo debe definirse «concomitante» (misma
sesión de vacunación u otra ventana de tiempo)?}%
{Concomitante significa otra vacuna administrada en la misma notificación; no se dispone de información
sobre otros tratamientos.}{3.4}

\Seccion{F. Cobertura, incidencia y denominadores (\Fichero{BD para publicación lotes final.xlsx})}
\Pregunta{cobertura}{Denominador de la hoja de cobertura}%
{Hoja «Cobertura de vacunación» (1986--2025): columna B sin encabezado con valores de 130\,233 a
280\,023 (15 valores) y columna «Cobertura» (9 años sin dato; mínimo 82,9).}%
{¿La columna B es la población objetivo, las dosis administradas de VA-MENGOC-BC o ambas según el año?
¿Se dispone de las dosis administradas de 2023, 2024 y 2025 (y de 2026 cuando concluya el año)? ¿El
valor mínimo de cobertura (82,9) corresponde a un año con campaña documentada?}%
{Se interpreta como dosis administradas; los años sin dato no tienen tasa de notificación y los años con
menos de 12 meses de notificaciones (2023, 2026) tampoco.}{3.6}

\Pregunta{colhoja}{Columnas sin encabezado y «Referencia» de la hoja de cobertura}%
{Hoja «Cobertura de vacunación»: columna «Referencia» (9 etiquetas) y columnas L, M y N sin encabezado.}%
{¿Qué contienen estas columnas? ¿Alguna debe incorporarse al análisis? ¿La columna N (un único valor
repetido en 7 filas) marca alguna condición administrativa?}%
{Se descartan y se informa como columnas sin reconocer.}{3.0}

\Pregunta{incidencia}{Serie de incidencia de enfermedad meningocócica}%
{Hoja 2 (incidencia y mortalidad; encabezado en la fila 5): filas de 1970 a 2024, aunque el título
indica «1980, 1985--2024»; columnas «Tasax100000 habitantes» (dos veces) y «Defunsiones».}%
{¿Las filas de 1970 a 1979 son datos reales o un artefacto de importación? ¿La segunda columna
«Tasax100000 habitantes» es la tasa de mortalidad? En «Defunsiones», ¿las celdas vacías son cero o
valores menores que 5 omitidos? ¿La fuente de la serie es la Dirección de Registros Médicos y
Estadísticas de Salud (año 2025)?}%
{Se usa toda la serie; la segunda tasa se asigna a mortalidad y las celdas vacías de defunciones se tratan
como no disponibles.}{3.8}

\Pregunta{tasaref}{Tasa de notificación de referencia del programa}%
{Resultados de tasas de notificación por 100\,000 dosis (primer artículo).}%
{El análisis compara la tasa de VA-MENGOC-BC con una referencia de 50 notificaciones por 100\,000 dosis,
tomada de Cruz-Rodríguez y cols. (2021) y referida a las vacunas infantiles en conjunto. ¿Confirma
este valor y su alcance, o existe una referencia nacional más adecuada y actual?}%
{Se usa 50 por 100\,000 dosis solo como referencia descriptiva.}{3.0}

\Seccion{G. Entregas mensuales y revisión de señales}
\Pregunta{entregas}{Formato y política de las entregas mensuales}%
{Libros de 2026 y sucesivos (\Fichero{EA 2026 TOTAL.xlsx}).}%
{¿Cada entrega mensual es un libro acumulado del año (con todos los meses desde enero) o contiene solo
el mes nuevo? ¿Se corrigen o amplían retroactivamente los meses ya entregados? ¿Mantendrá el
nombre «EA AAAA TOTAL.xlsx» y la hoja y las columnas del formulario de 2025? ¿Cuál es la fecha de
corte de cada entrega?}%
{Se espera un libro por año, que sustituye al anterior (no se combinan versiones acumuladas); la
ventana de estudio sigue a los datos y los años incompletos se señalan como parciales.}{3.8}

\Pregunta{comite}{Revisión de las señales por el comité nacional}%
{Resultados de desproporcionalidad del primer artículo (eventos no descritos como esperados).}%
{¿Qué comité o instancia revisará caso a caso, con las historias clínicas, las señales que no
correspondan a reacciones esperadas, y en qué plazo? ¿Debe incluirse su conclusión en el artículo?}%
{Las señales se presentan como hipótesis para revisión clínica; el análisis no incluye revisión caso a
caso.}{3.0}

\Seccion{Observaciones generales}
\noindent
\TextField[name=observaciones,multiline=true,width=\linewidth,height=5cm,charsize=10pt,%
  bordercolor=gris,backgroundcolor=fondo,align=0]{}\par\bigskip

Agradecemos su colaboración. Las respuestas permitirán completar los manuscritos con datos
verificados y documentar cada decisión metodológica.

\medskip
\textit{Elaborado por:} Richard Alejandro Matos Arderí y Abel Ponce González, Instituto Finlay de
Vacunas, DICEI.
\end{Form}
\end{document}
